\documentclass[12pt]{article}
\usepackage[a4paper,margin=2.5cm]{geometry}
\usepackage{graphicx}
\usepackage{amsmath,amssymb}
\usepackage{hyperref}

\title{Identifikace strukturních motivů chemických látek klíčových pro rozhodování u regresních QSPR a QSAR modelů}
\author{Michael Kolakovský}
\date{\today}

\begin{document}
\maketitle
\thispagestyle{empty} % optional: no page number on title page
\newpage

\begin{abstract}
  Quantitative Structure–Property Relationship (QSPR) a Quantitative Structure–Activity Relationship regresní model jsou často používané pro předpovídání fyzikálních vlastností a biologické aktivity chemických látek.
  Tyto modely se například používají pro farmaceutické účely, kde zjištění těchto vlastností by mohlo stát velký obnos peněz, tyto modely mohou docela je odhadnout.
  Modely dosahují vysoké přesnosti, jak dokazuje část této práce, ale jejich rozhodování je často složité pochopit.
  Cílem této práce je identifikovat klíčové strukturní motivy chemických látek pomocí metodiky SHAP a Boruta-SHAP.
\end{abstract}
\newpage
  
\begin{Použité datové sady}
  Pro QSPR modely byla využita datová sada AqSolDB, která obsahuje přibližně 10 000 chemických látek s jejich vlastností rozpustností ve vodě, která se dá použít například pro farmaceutické účely.
\end{Použité datové sady}

\begin{Random Forest}
  Pro jeden z prvních modelů byla použita metoda Random Forest.
  Pro QSPR byla užita již zmíněná datová sada AqSolDB.

\end{Random Forest}
  

\end{document}
